\chapter{Determinants as Volume}\label{linear:ch:volume}

\begin{center}\itshape
``Then you should say what you mean,'' the March Hare went on.\\
``I do,'' Alice hastily replied; ``at least---at least I mean what I say---that's the same thing, you know.''
\upshape(Lewis Carroll, \emph{Alice's Adventures in Wonderland}, 1865)
\end{center}

A story used to circulate in Victorian drawing rooms. Queen Victoria, so it went, was so delighted by \emph{Alice's
Adventures in Wonderland} that she asked for the author's next book to be sent to her, and received, a few months later,
\emph{An Elementary Treatise on Determinants}. The author, the Reverend Charles Lutwidge Dodgson, lecturer in mathematics
at Christ Church, Oxford, denied the story firmly in a later edition of one of his books. But the treatise was real,
and so was a short paper he read to the Royal Society in 1866, \emph{Condensation of Determinants}, describing a way to
compute an $n\times n$ determinant by repeatedly replacing each $2\times2$ block of adjacent entries by its determinant
and dividing by an entry from two steps before. It looks like a party trick. It is exact, and it was forgotten for a
century.

In the early 1980s, David Robbins and Howard Rumsey at the Institute for Defense Analyses were playing with Dodgson's
condensation, and asked what happens if the $2\times2$ determinant $ad-bc$ is replaced by $ad+\lambda bc$. The answer
was a sum over strange matrices of $0$s, $1$s and $-1$s, in which the non-zero entries alternate in sign along every row
and column. They counted these \emph{alternating sign matrices}: $1,2,7,42,429,7436,\dots$ With William Mills they
guessed a formula, and it took thirteen years and an army of referees before Doron Zeilberger proved it in 1996. The
same years saw the discovery that random tilings of regions by lozenges and dominoes have frozen corners and a
disordered core separated by a sharp curve, the \emph{arctic circle}. Determinants were behind all of it.

This chapter rebuilds the determinant from its geometric meaning, the signed volume of a parallelepiped, and follows it
to the counting of paths that never cross. We define the exterior powers of a space, in which the determinant becomes
the simplest possible object, a number by which a map scales the top power; the theorem of Cauchy and Binet then becomes
the statement that this construction respects composition. We prove the identities of Jacobi and Desnanot behind
Dodgson's condensation. And we prove the lemma of Lindstr\"om, Gessel and Viennot, which says that a determinant of path
counts counts families of paths that do not meet, and use it to count lozenge tilings.

\begin{think}
A hexagon with all sides of length $2$ and all angles $120^\circ$ is cut into $24$ small equilateral triangles. In how
many ways can it be tiled by \emph{lozenges}, the rhombi made of two adjacent small triangles? Draw some tilings; you
will find that each looks like a pile of cubes in a corner. The answer is $20$. For the hexagon of side $3$ it is $980$,
and for side $10$ it is a $34$-digit number. Section~\ref{linear:sec:lgv} computes all of them as determinants.
\end{think}

\section{Volume and orientation}

Let $v_1,\dots,v_n\in\R^n$. The \emph{parallelepiped} they span is $\{\sum_it_iv_i:0\le t_i\le1\}$. Its volume should
satisfy three rules, which are already visible in the plane (Figure~\ref{linear:fig:shear}).
\begin{enumerate}[label=(\roman*)]
\item Stretching one edge by a factor $\lambda\ge0$ multiplies the volume by $\lambda$.
\item Adding to one edge a multiple of another, a \emph{shear}, does not change the volume: the slices parallel to the
other edges are merely slid along, as in Cavalieri's principle.
\item The unit cube has volume $1$.
\end{enumerate}
Rule (i) with a negative $\lambda$ would contradict positivity. The fix, which turns out to be enormously convenient, is
to allow negative volumes: an ordered family $(v_1,\dots,v_n)$ has an \emph{orientation}, and a signed volume
$D(v_1,\dots,v_n)$ that is multiplied by $\lambda$ for every real $\lambda$ in rule (i). Then (i) and (ii) say that $D$
is linear in each edge (we take this as the axiom that rules (i) and (ii) suggest) and that $D$ vanishes when two
edges coincide (a shear makes one of them zero, and a
parallelepiped with a zero edge is flat). By the uniqueness theorem of Volume~0 (Theorem~7.2), the only such function with
$D(e_1,\dots,e_n)=1$ is the determinant:
\[
D(v_1,\dots,v_n)=\det(v_1,\dots,v_n).
\]

\begin{figure}[ht]
\centering
\fig{linear-ch02-shear}
\caption{Replacing $v_2$ by $v_2+v_1$ slides the parallelogram along its base: same base, same height, same area. In
coordinates the area is $|\det(v_1,v_2)|$.}
\label{linear:fig:shear}
\end{figure}

\begin{remark}
That the Lebesgue measure of the image of a box under a linear map $T$ is $|\det T|$ times its volume, which turns the
heuristic above into a theorem, is proved in Volume~6, Chapter~8; the change of variables formula for integrals rests on
it. In this chapter we only use the algebra of determinants, and the word ``volume'' serves as a guide.
\end{remark}

The sign of $\det(v_1,\dots,v_n)$ is the \emph{orientation} of the basis: in the plane, positive when turning from $v_1$
to $v_2$ is counterclockwise. Here is the first dividend of signed volumes.

\begin{example}[the shoelace formula]\label{linear:ex:shoelace}
Let $P$ be a polygon with vertices $p_1,\dots,p_m$ listed counterclockwise, $p_{m+1}=p_1$, and $p_k=(x_k,y_k)$. Then
\[
\operatorname{area}(P)=\frac12\sum_{k=1}^m\det(p_k,p_{k+1})=\frac12\sum_{k=1}^m(x_ky_{k+1}-x_{k+1}y_k).
\]
Indeed $\frac12\det(p_k,p_{k+1})$ is the signed area of the triangle $0p_kp_{k+1}$. For a convex polygon with the origin
inside, the triangles tile $P$ with positive signs. In general, the signed areas of the triangles with apex $0$ add up
to the area of $P$ because every point outside $P$ is covered as many times positively as negatively (the winding number
of the boundary around it is $0$), and every point inside is covered once more positively than negatively. The formula
needs no triangulation, works for any simple polygon, and is how a computer measures the area of a field from the GPS
coordinates of its corners.
\end{example}

\section{Exterior powers}\label{linear:sec:exterior}

The determinant measures $n$-dimensional volume in an $n$-dimensional space. What measures the area of a parallelogram
in $\R^3$, or more generally the $k$-dimensional size of $k$ vectors in an $n$-dimensional space? The answer is not a
number but a vector in a new space, and the construction of that space is the purpose of this section.

Let $V$ have a basis $e_1,\dots,e_n$, and let $1\le k\le n$. For a $k$-element subset $I=\{i_1<\dots<i_k\}$ of
$[n]=\{1,\dots,n\}$ we use a symbol $e_I$, and we let $\Lambda^kV$ be the space with basis the $\binom nk$ symbols $e_I$.
For vectors $v_1,\dots,v_k\in V$, let $M$ be the $n\times k$ matrix of their coordinates and $M_I$ the $k\times k$
submatrix formed by the rows in $I$; we define their \emph{exterior product}
\[
v_1\wedge\dots\wedge v_k=\sum_{|I|=k}\det(M_I)\,e_I\in\Lambda^kV.
\]
Its coordinates are the $k\times k$ minors of $M$, the signed $k$-volumes of the projections of the parallelepiped onto
the coordinate $k$-planes. For $k=n$ there is one minor, and $v_1\wedge\dots\wedge v_n=\det(v_1,\dots,v_n)\,e_{[n]}$; for
$k=1$, $v_1$ is $v_1$. For $k=2$ and $n=3$, the coordinates of $v\wedge w$ are $v_1w_2-v_2w_1$, $v_1w_3-v_3w_1$ and
$v_2w_3-v_3w_2$: up to order and sign, the cross product $v\times w$, whose length is the area of the parallelogram
(Exercise~\ref{linear:ex:lagrangeid}).

\begin{theorem}[exterior powers]\label{linear:thm:exterior}
\begin{enumerate}[label=(\alph*)]
\item The map sending $(v_1,\dots,v_k)$ to $v_1\wedge\dots\wedge v_k$ is multilinear and alternating, and
$e_{i_1}\wedge\dots\wedge e_{i_k}=e_I$ for $i_1<\dots<i_k$.
\item (Universal property.) For every multilinear alternating map $\varphi:V^k\to U$ there is a unique linear map
$\tilde\varphi:\Lambda^kV\to U$ with $\varphi(v_1,\dots,v_k)=\tilde\varphi(v_1\wedge\dots\wedge v_k)$.
\item The vectors $v_1,\dots,v_k$ are linearly independent if and only if $v_1\wedge\dots\wedge v_k\ne0$.
\end{enumerate}
\end{theorem}
\begin{proof}
(a) Each coordinate $\det(M_I)$ is multilinear and alternating in the columns of $M$, which are the $v_j$. For
$v_j=e_{i_j}$, $M_J$ is the identity when $J=I$ and has a zero row otherwise.

(b) Uniqueness: the $e_I=e_{i_1}\wedge\dots\wedge e_{i_k}$ form a basis, and $\tilde\varphi(e_I)$ must be
$\varphi(e_{i_1},\dots,e_{i_k})$. Existence: define $\tilde\varphi$ by this formula on the basis. To check
$\varphi(v_1,\dots,v_k)=\tilde\varphi(v_1\wedge\dots\wedge v_k)$, expand each $v_j=\sum_im_{ij}e_i$ by multilinearity:
$\varphi(v_1,\dots,v_k)=\sum_fm_{f(1)1}\cdots m_{f(k)k}\,\varphi(e_{f(1)},\dots,e_{f(k)})$ over all maps
$f:[k]\to[n]$. Non-injective $f$ contribute $0$ by alternation. An injective $f$ has image a set $I$ and is
$f(j)=i_{\sigma(j)}$ for a unique permutation $\sigma$ of $[k]$; since an alternating map changes sign under a
transposition of two arguments, $\varphi(e_{f(1)},\dots,e_{f(k)})=\varepsilon(\sigma)\varphi(e_{i_1},\dots,e_{i_k})$. Grouping the
terms by $I$, we get $\sum_I\bigl(\sum_\sigma\varepsilon(\sigma)\prod_jm_{i_{\sigma(j)}j}\bigr)\varphi(e_{i_1},\dots,e_{i_k})=\sum_I\det(M_I)\tilde\varphi(e_I)$,
as required.

(c) If the $v_j$ are dependent, one is a combination of the others and multilinearity plus alternation give
$v_1\wedge\dots\wedge v_k=0$. If they are independent, complete them into a basis $v_1,\dots,v_n$ of $V$; the matrix of all
$n$ vectors is invertible, and if all $k\times k$ minors of its first $k$ columns vanished, the Laplace expansion along
those columns (Exercise~\ref{linear:ex:laplace}) would give determinant $0$. So some $\det(M_I)\ne0$.
\end{proof}

The universal property says that $\Lambda^kV$ is ``the'' space through which alternating $k$-linear maps factor. In
particular it does not depend on the basis we started from, up to a canonical isomorphism: two spaces with the universal
property are isomorphic by the maps the property provides. The key consequence is functoriality.

\begin{definition}\label{linear:def:compound}
Let $f:V\to W$ be linear. The \emph{$k$-th exterior power} of $f$ is the linear map $\Lambda^kf:\Lambda^kV\to\Lambda^kW$
with $\Lambda^kf(v_1\wedge\dots\wedge v_k)=f(v_1)\wedge\dots\wedge f(v_k)$; it exists and is unique by
Theorem~\ref{linear:thm:exterior}(b) applied to the alternating map $(v_j)\mapsto f(v_1)\wedge\dots\wedge f(v_k)$. For a matrix
$A\in M_{m,n}(\mathbb K)$, the \emph{$k$-th compound matrix} $C_k(A)$ is the $\binom mk\times\binom nk$ matrix of its
$k\times k$ minors: $C_k(A)_{I,J}=\det A_{I,J}$ (rows $I$, columns $J$, both listed in increasing order).
\end{definition}

\begin{proposition}\label{linear:prop:compound}
If $f$ has matrix $A$ in bases $(e_j)$ and $(u_i)$, then $\Lambda^kf$ has matrix $C_k(A)$ in the bases $(e_J)$, $(u_I)$.
Moreover $\Lambda^k(g\circ f)=\Lambda^kg\circ\Lambda^kf$ and $\Lambda^k\mathrm{id}=\mathrm{id}$. For $V=W$ of dimension $n$,
$\Lambda^nf$ is multiplication by $\det f$ on the line $\Lambda^nV$.
\end{proposition}
\begin{proof}
$\Lambda^kf(e_J)=f(e_{j_1})\wedge\dots\wedge f(e_{j_k})$; the coordinate matrix of these $k$ vectors consists of the columns
$J$ of $A$, so by definition of the exterior product the $u_I$-coordinate is $\det A_{I,J}$. The composition rule holds on
the generators $v_1\wedge\dots\wedge v_k$, which span $\Lambda^kV$. For $k=n$, $C_n(A)$ is the $1\times1$ matrix $(\det A)$.
\end{proof}

The last statement is the basis-free definition of the determinant: on the line $\Lambda^nV$, every linear map acts by a
scalar, and that scalar is $\det f$. Multiplicativity $\det(gf)=\det g\det f$ is now a special case of functoriality.
The general case is a famous formula, found by Binet and Cauchy independently in 1812.

\begin{theorem}[Cauchy--Binet]\label{linear:thm:cauchybinet}
Let $A\in M_{m,n}(\mathbb K)$ and $B\in M_{n,p}(\mathbb K)$. Then $C_k(AB)=C_k(A)C_k(B)$; that is, for $|I|=|J|=k$,
\[
\det(AB)_{I,J}=\sum_{|S|=k}\det A_{I,S}\,\det B_{S,J}.
\]
In particular, for $A\in M_{k,n}$ and $B\in M_{n,k}$,
\[
\det(AB)=\sum_{|S|=k}\det A_{[k],S}\det B_{S,[k]},
\]
which is $0$ when $n<k$.
\end{theorem}
\begin{proof}
Proposition~\ref{linear:prop:compound} applied to the maps $\mathbb K^p\xrightarrow{B}\mathbb K^n\xrightarrow{A}\mathbb K^m$:
the matrix of $\Lambda^k(AB)$ is $C_k(AB)$ and that of $\Lambda^kA\circ\Lambda^kB$ is $C_k(A)C_k(B)$. When $n<k$, $\Lambda^k\mathbb K^n=0$.
\end{proof}

\begin{corollary}[Pythagoras for $k$-volumes]\label{linear:cor:pythagoras}
For $M\in M_{n,k}(\R)$, $\det(M^{\mathsf T}M)=\sum_{|I|=k}\det(M_I)^2$. In words: the square of the $k$-volume of the
parallelepiped spanned by the columns of $M$ is the sum of the squares of the $k$-volumes of its projections on the
coordinate $k$-planes.
\end{corollary}
\begin{proof}
Cauchy--Binet with $A=M^{\mathsf T}$, $B=M$. That $\sqrt{\det(M^{\mathsf T}M)}$ is the $k$-volume is proved in Appendix~D
(Gram determinants). For $k=1$ this is Pythagoras' theorem; for $k=2$, $n=3$ it says that the squared area of a
triangle in space is the sum of the squared areas of its three shadows on the coordinate planes, a theorem that de Gua
stated in 1783 for the right-angled corner of a tetrahedron.
\end{proof}

\section{Minors of the inverse and Dodgson's condensation}\label{linear:sec:condensation}

For a square matrix $A$ and sets $I,J\subseteq[n]$ of the same size, write $A_{I,J}$ for the submatrix with rows $I$ and
columns $J$, $I^c=[n]\setminus I$, and $s(I)=\sum_{i\in I}i$. The adjugate formula $A^{-1}=\operatorname{adj}(A)/\det A$
expresses each entry of $A^{-1}$ through a minor of size $n-1$ of $A$. Jacobi generalised it in 1833 to minors of any size.

\begin{theorem}[Jacobi's complementary minors]\label{linear:thm:jacobi}
Let $A\in GL_n(\mathbb K)$ and $I,J\subseteq[n]$ with $|I|=|J|=k$. Then
\[
\det(A^{-1})_{J,I}=(-1)^{s(I)+s(J)}\,\frac{\det A_{I^c,J^c}}{\det A}.
\]
\end{theorem}
\begin{proof}
First let $I=J=[k]$. Write $A$ and $A^{-1}$ in blocks of sizes $k$ and $n-k$:
$A=\begin{pmatrix}P&Q\\R&S\end{pmatrix}$, $A^{-1}=\begin{pmatrix}W&X\\Y&Z\end{pmatrix}$. From $AA^{-1}=I$ we get $PW+QY=I_k$ and
$RW+SY=0$, hence
\[
A\begin{pmatrix}W&0\\Y&I_{n-k}\end{pmatrix}=\begin{pmatrix}I_k&Q\\0&S\end{pmatrix}.
\]
Both block matrices are block triangular; taking determinants, $\det A\cdot\det W=\det S$, which is the claim. In
general, let $\sigma$ be the permutation that lists the elements of $I$ in increasing order followed by those of $I^c$
in increasing order, and $\tau$ the same for $J$; let $P_\sigma,P_\tau$ be the corresponding permutation matrices, chosen
so that $B=P_\sigma AP_\tau^{-1}$ has $B_{[k],[k]}=A_{I,J}$ and $B_{[k]^c,[k]^c}=A_{I^c,J^c}$. Then
$B^{-1}=P_\tau A^{-1}P_\sigma^{-1}$ has $(B^{-1})_{[k],[k]}=(A^{-1})_{J,I}$, and $\det B=\varepsilon(\sigma)\varepsilon(\tau)\det A$. The first
case gives $\det(A^{-1})_{J,I}=\varepsilon(\sigma)\varepsilon(\tau)\det A_{I^c,J^c}/\det A$. Finally, $\sigma$ is a product of
$\sum_r(i_r-r)=s(I)-k(k+1)/2$ adjacent transpositions (move $i_1$ to position $1$, then $i_2$ to position $2$, and so on),
so $\varepsilon(\sigma)\varepsilon(\tau)=(-1)^{s(I)+s(J)-k(k+1)}=(-1)^{s(I)+s(J)}$.
\end{proof}

Write $A^{i}_{j}$ for $A$ with row $i$ and column $j$ deleted, and $A^{1,n}_{1,n}$ for $A$ with rows and columns $1$
and $n$ deleted. For $k=1$, with $I=\{i\}$ and $J=\{j\}$, the theorem says $(A^{-1})_{ji}=(-1)^{i+j}\det A^{i}_{j}/\det A$:
this is the \emph{adjugate formula} $A^{-1}=\operatorname{adj}(A)/\det A$, where $\operatorname{adj}(A)$ is the matrix with
$(j,i)$ entry $(-1)^{i+j}\det A^{i}_{j}$. It is worth recording the version that needs no inverse.

\begin{corollary}[adjugate identity]\label{linear:cor:adjugate}
For every $A\in M_n(R)$, $R$ any commutative ring, $A\operatorname{adj}(A)=\operatorname{adj}(A)A=(\det A)I$.
\end{corollary}
\begin{proof}
Each entry of $A\operatorname{adj}(A)-(\det A)I$ is a polynomial $p$ with integer coefficients in the $n^2$ entries of $A$.
For $A\in GL_n(\Q)$ the entry vanishes by the case $k=1$ of Theorem~\ref{linear:thm:jacobi}, so the polynomial $p\cdot\det$
vanishes at every point of $\Q^{n^2}$, hence is the zero polynomial (a non-zero polynomial over an infinite field has a
non-root, by induction on the number of variables), hence $p=0$ since $\Z[a_{11},\dots,a_{nn}]$ is an integral domain.
A polynomial identity with integer coefficients that holds in $\Z[a_{11},\dots,a_{nn}]$ holds after any substitution
of the variables in any commutative ring. The same argument gives $\operatorname{adj}(A)A=(\det A)I$.
\end{proof}

Chapter~\ref{linear:ch:eigen} uses the corollary over the polynomial ring $\mathbb K[X]$ to prove Cayley--Hamilton. For $k=2$ and
the corners of the matrix, Theorem~\ref{linear:thm:jacobi} becomes an identity that Desnanot found in 1819 for $n\le6$ and
Jacobi proved in general.

\begin{theorem}[Desnanot--Jacobi identity]\label{linear:thm:desnanot}
For every $A\in M_n(\mathbb K)$, $n\ge2$ (with $\det A^{1,n}_{1,n}=1$ when $n=2$),
\[
\det A\cdot\det A^{1,n}_{1,n}=\det A^1_1\det A^n_n-\det A^1_n\det A^n_1.
\]
\end{theorem}
\begin{proof}
Suppose first $A$ invertible, and apply Theorem~\ref{linear:thm:jacobi} with $I=J=\{1,n\}$: the sign is $(-1)^{2(1+n)}=1$, so the
$2\times2$ determinant of $A^{-1}$ on rows and columns $\{1,n\}$ is $\det A^{1,n}_{1,n}/\det A$. By the adjugate formula
$(A^{-1})_{ij}=(-1)^{i+j}\det A^j_i/\det A$, so that determinant is
\[
\frac{\det A^1_1\det A^n_n-(-1)^{2(1+n)}\det A^n_1\det A^1_n}{(\det A)^2}.
\]
Equating the two expressions and multiplying by $(\det A)^2$ gives the identity. Both sides are polynomials in the
entries of $A$. For an arbitrary $A$, consider $A+tI$ for an indeterminate $t$: both sides become polynomials in $t$, and
they agree for every $t\in\mathbb K$ with $\det(A+tI)\ne0$, that is, for all but at most $n$ values. If $\mathbb K$ is
infinite, two polynomials agreeing at infinitely many points are equal, and $t=0$ gives the identity for $A$. For a finite
field, do the same computation over the infinite field $\mathbb K(t)$ of rational functions, where $A+tI$ is invertible
because $\det(A+tI)$ is a monic polynomial of degree $n$ in $t$; then set $t=0$.
\end{proof}

\begin{corollary}[Dodgson's condensation]\label{linear:cor:dodgson}
For $A\in M_n(\mathbb K)$ and $0\le k\le n$, let $A^{(k)}$ be the $(n-k+1)\times(n-k+1)$ matrix of the \emph{contiguous}
$k\times k$ minors of $A$: $A^{(k)}_{ij}$ is the determinant of the block of $A$ with rows $i,\dots,i+k-1$ and columns
$j,\dots,j+k-1$. So $A^{(1)}=A$, $A^{(0)}$ is the all-ones matrix of size $n+1$, and $A^{(n)}=(\det A)$. Then
\[
A^{(k+1)}_{ij}\,A^{(k-1)}_{i+1,j+1}=A^{(k)}_{ij}A^{(k)}_{i+1,j+1}-A^{(k)}_{i,j+1}A^{(k)}_{i+1,j}.
\]
\end{corollary}
\begin{proof}
Apply Theorem~\ref{linear:thm:desnanot} to the contiguous $(k+1)\times(k+1)$ block with top-left corner $(i,j)$: its four corner
minors of size $k$ are the four $A^{(k)}$ entries on the right, and its interior is the block counted by
$A^{(k-1)}_{i+1,j+1}$.
\end{proof}

Dodgson's algorithm computes $A^{(2)},A^{(3)},\dots$ in turn, each entry from four entries of the previous matrix
divided by one entry of the matrix before. It uses only $2\times2$ determinants and is easy to run by hand, or in
parallel. Its weakness is division: if an interior entry vanishes the algorithm stops, although the determinant may
be perfectly fine (Exercise~\ref{linear:ex:dodgson}). Over the integers, every division is exact, since the quotient is a
minor.

\begin{remark}
Robbins and Rumsey's $\lambda$-determinant replaces the right-hand side of the condensation rule by
$A^{(k)}_{ij}A^{(k)}_{i+1,j+1}+\lambda A^{(k)}_{i,j+1}A^{(k)}_{i+1,j}$. They proved that the result is a sum over
alternating sign matrices, and Mills, Robbins and Rumsey conjectured in 1983 that there are exactly
$\prod_{j=0}^{n-1}\frac{(3j+1)!}{(n+j)!}$ alternating sign matrices of size $n$. Zeilberger's proof (1996) runs to $84$
pages; Kuperberg found a shorter one the same year through a model of statistical physics, the six-vertex ice model.
Neither is within the scope of this book; see the book by Bressoud listed at the end of the chapter.
\end{remark}

\section{Paths that never meet}\label{linear:sec:lgv}

Let $G$ be a finite directed graph without directed cycles, with a weight $w(e)$ in a commutative ring on each edge
(take $w=1$ to count). The weight of a directed path is the product of the weights of its edges, and for vertices $a,b$
let
\[
e(a,b)=\sum_{\text{paths }P:a\to b}w(P),
\]
with $e(a,a)=1$ (the path of length $0$). Given \emph{sources} $A_1,\dots,A_k$ and \emph{sinks} $B_1,\dots,B_k$, a
\emph{path system} is a permutation $\sigma$ of $[k]$ together with paths $P_i:A_i\to B_{\sigma(i)}$; it is
\emph{non-intersecting} if no two of its paths share a vertex. Its weight is $w(P_1)\cdots w(P_k)$ and its sign is
$\varepsilon(\sigma)$.

\begin{theorem}[Lindstr\"om, 1973; Gessel--Viennot, 1985]\label{linear:thm:lgv}
With $M=(e(A_i,B_j))_{i,j}$,
\[
\det M=\sum_{\text{non-intersecting systems }(\sigma,P)}\varepsilon(\sigma)\,w(P_1)\cdots w(P_k).
\]
\end{theorem}
\begin{proof}
Expanding the determinant and each entry,
\[
\det M=\sum_\sigma\varepsilon(\sigma)\prod_ie(A_i,B_{\sigma(i)})=\sum_{(\sigma,P)}\varepsilon(\sigma)\,w(P_1)\cdots w(P_k),
\]
the sum over \emph{all} path systems. We build an involution $\Phi$ on the set of intersecting systems that preserves
the weight and reverses the sign; then the intersecting systems cancel in pairs and only the non-intersecting ones
remain. Given an intersecting system, let $i_0$ be the smallest index such that $P_{i_0}$ meets another path; let $v$
be the first vertex on $P_{i_0}$ that lies on another path; and let $j_0$ be the smallest index $j\ne i_0$ such that
$P_j$ passes through $v$. Now swap the tails: replace $P_{i_0}$ by the part of $P_{i_0}$ up to $v$ followed by the part
of $P_{j_0}$ after $v$, and $P_{j_0}$ by the part of $P_{j_0}$ up to $v$ followed by the part of $P_{i_0}$ after $v$
(Figure~\ref{linear:fig:lgv}). The new paths go from $A_{i_0}$ to $B_{\sigma(j_0)}$ and from $A_{j_0}$ to $B_{\sigma(i_0)}$, so the new
permutation is $\sigma\circ(i_0\,j_0)$, of opposite sign. The multiset of edges used is the same, so the weight is the
same. (Since $G$ has no directed cycles, a path visits $v$ only once, and the new walks are again paths.)

It remains to see that $\Phi$ is an involution: applied to the new system, it must choose the same $i_0$, $v$ and
$j_0$. The set of vertices covered by the system, and which of them lie on two or more paths, are unchanged. Path
$i_0$ still meets another path; a path $P_i$ with $i<i_0$ is unchanged and met no other path before, and still meets
none, because the vertices covered by the union of the other paths are unchanged as a set. So $i_0$ is chosen again.
The initial segment of $P_{i_0}$ up to $v$ is unchanged, and none of its vertices before $v$ lay on another path; this is
still true. So $v$ is chosen again. Finally the paths through $v$ are the same indices as before, so $j_0$ is chosen
again, and swapping the tails again restores the original system.
\end{proof}

\begin{figure}[ht]
\centering
\fig{linear-ch02-lgv}
\caption{The tail swap. Two paths meet at the marked vertex; exchanging what comes after it changes the permutation by a
transposition and keeps the edges, so the two systems cancel in the determinant.}
\label{linear:fig:lgv}
\end{figure}

The theorem is most useful when only the identity permutation admits non-intersecting systems, for then $\det M$ simply
\emph{counts} them (with weights). This happens in the plane when the sources and the sinks are suitably ordered. We
state it for the lattice $\Z^2$ with unit steps east $(1,0)$ and north $(0,1)$, the setting of every example below.

\begin{lemma}\label{linear:lem:planar}
Let $A_1,\dots,A_k$ and $B_1,\dots,B_k$ be points of $\Z^2$ such that $A_{i+1}$ is weakly north-west of $A_i$ (smaller or
equal $x$, larger or equal $y$, not both equal), and likewise $B_{i+1}$ of $B_i$. If $i<j$, every north-east path from
$A_i$ to $B_l$ meets every north-east path from $A_j$ to $B_m$ whenever $l>m$. Consequently, in a non-intersecting system
of such paths, $\sigma=\mathrm{id}$.
\end{lemma}
\begin{proof}
A north-east path visits exactly one point on each anti-diagonal $x+y=s$ between the values of $x+y$ at its two ends;
write $x_P(s)$ for the $x$-coordinate of that point. Along the path, $x_P(s+1)-x_P(s)\in\{0,1\}$, and $y=s-x_P(s)$ is
non-decreasing too. Let $P$ go from $A_i$ to $B_l$ and $Q$ from $A_j$ to $B_m$, with $i<j$ and $l>m$, so that $A_j$ is weakly
north-west of $A_i$ and $B_l$ weakly north-west of $B_m$. Let $s_0=\max(s(A_i),s(A_j))$ and $s_1=\min(s(B_l),s(B_m))$, where
$s(p)$ is the sum of the coordinates of $p$.

The two paths both visit the anti-diagonals from $s_0$ to $s_1$: indeed $x(B_m)\ge x(B_l)\ge x(A_i)$ and $y(B_m)\ge y(A_j)\ge
y(A_i)$ give $s(B_m)\ge s(A_i)$, and similarly $s(B_l)\ge s(A_j)$, so $s_0\le s_1$. Let $\delta(s)=x_P(s)-x_Q(s)$ for
$s_0\le s\le s_1$. At $s_0$: if $s_0=s(A_j)$, then $x_Q(s_0)=x(A_j)\le x(A_i)\le x_P(s_0)$; if $s_0=s(A_i)>s(A_j)$, then
$x_P(s_0)=x(A_i)$ and $x_Q(s_0)=s_0-y_Q\le s_0-y(A_j)\le s_0-y(A_i)=x(A_i)$. Either way $\delta(s_0)\ge0$. The symmetric argument at
the other end gives $\delta(s_1)\le0$. Since $\delta$ changes by at most $1$ at each step, it vanishes at some $s$, where
the two paths visit the same point. For the consequence: if $\sigma\ne\mathrm{id}$, there are $i<j$ with $\sigma(i)>\sigma(j)$, and
the paths from $A_i$ and $A_j$ meet.
\end{proof}

The number of north-east paths from $(x,y)$ to $(x',y')$ is $\binom{(x'-x)+(y'-y)}{x'-x}$: choose which of the steps go
east. (It is $0$ if $x'<x$ or $y'<y$, with the convention $\binom nk=0$ for $k<0$ or $k>n$.) So the LGV lemma turns
families of non-crossing lattice paths into determinants of binomial coefficients.

\begin{example}\label{linear:ex:binomialdet}
The determinant $\det\bigl(\binom{i+j}{i}\bigr)_{0\le i,j\le n}$ equals $1$. Take $A_i=(-i,0)$ and $B_j=(0,j)$: a path from
$A_i$ to $B_j$ has $i$ east steps and $j$ north steps, so $e(A_i,B_j)=\binom{i+j}i$. The sources move west and the sinks
move north, so Lemma~\ref{linear:lem:planar} applies and non-intersecting systems send $A_i$ to $B_i$. There is exactly one.
The path from $A_0=B_0=(0,0)$ is the empty path. Suppose the paths from $A_0,\dots,A_{i-1}$ are the \emph{hooks} $H_k$
that go north from $(-k,0)$ to $(-k,k)$, then east to $(0,k)$. The path from $A_i=(-i,0)$ must leave the column $x=-i$ by
an east step from some $(-i,y)$ to $(-i+1,y)$; the points $(-i+1,y)$ with $y\le i-1$ lie on $H_{i-1}$, so $y\ge i$, and as the
path ends at height $i$, $y=i$. So the path is the hook $H_i$, which meets no earlier hook (they live at heights below
$i$ and in columns east of $-i$). The determinant counts this single system: it is $1$.
\end{example}

\subsection*{Lozenge tilings}

We can now answer the opening question. Fix positive integers $a,b,c$ and let $H(a,b,c)$ be the hexagon with sides
$a,b,c,a,b,c$ (in this cyclic order) and all angles $120^\circ$, drawn on the triangular lattice. Two of its sides, of
length $c$ say, are vertical. A \emph{lozenge} is the union of two lattice triangles sharing an edge. There are three
kinds of lozenges, one for each direction of the shared edge, and exactly two kinds have a pair of vertical edges.

\begin{figure}[ht]
\centering
\fig{linear-ch02-lozenge}
\caption{A lozenge tiling of $H(4,5,3)$, seen as a pile of cubes in a $4\times5\times3$ box. The lozenges with vertical
edges (shaded) form three chains, one from each unit edge of the left vertical side to the right vertical side; their
midlines are drawn. The chains never meet, and they determine the tiling.}
\label{linear:fig:lozenge}
\end{figure}

\begin{theorem}[counting lozenge tilings]\label{linear:thm:lozenge}
The number of lozenge tilings of $H(a,b,c)$ is
\[
T(a,b,c)=\det\left(\binom{a+b}{a+i-j}\right)_{1\le i,j\le c}.
\]
\end{theorem}
\begin{proof}
\emph{Chains.} Take a tiling. Every lattice triangle has exactly one vertical edge. Call a lozenge \emph{vertical} if
it has vertical edges; it then has exactly two, one in each of its triangles. Start at a unit edge $\epsilon$ of the left
vertical side of the hexagon: the triangle inside the hexagon adjacent to $\epsilon$ belongs to some lozenge, which is
vertical because its triangle's vertical edge $\epsilon$ is not the shared edge (it is on the boundary). Leave that
lozenge through its other vertical edge, enter the next lozenge, which is again vertical for the same reason, and so on.
Each step moves strictly to the right, so the process ends on the right vertical side, the only other vertical part of
the boundary. This gives $c$ \emph{chains}. Two chains cannot share a lozenge, since a lozenge has only two vertical
edges and each chain enters and leaves through them. Conversely, every vertical lozenge lies on a chain: follow its
vertical edges to the left and to the right. And the non-vertical lozenges are determined by the chains: a triangle
$T$ not on a chain cannot have its vertical edge on the boundary (the boundary's vertical edges all start or end chains,
whose first and last lozenges contain the adjacent triangles), so the triangle $T'$ across that edge exists; $T'$ is not on a
chain either (a chain through $T'$ would cross the vertical edge of $T'$, which is the edge shared with $T$, and so would
contain $T$); and $T\cup T'$ must be the lozenge of the tiling containing $T$, since the lozenge containing $T$ is not
vertical, so it is glued along the vertical edge. So a tiling is the same as a family of $c$ disjoint chains, with the
remaining triangles paired across their vertical edges.

\emph{Paths.} Record a chain by the midpoints of the vertical edges it crosses. From one midpoint to the next, the chain
moves by one of two unit vectors: $v$, pointing right and down at $30^\circ$ below the horizontal, or $u$, pointing right
and up at $30^\circ$ above it. Call a $v$-step ``east'' and a $u$-step ``north''. The upward vertical unit vector is
$d=u-v$, so with the origin at a suitable midpoint, the midpoint $p+xv+yu$ gets coordinates $(x,y)\in\Z^2$, $d$ becomes
$(-1,1)$, and a chain becomes a north-east lattice path. Number the unit edges of the left side from bottom to top; their
midpoints are $A_i=(-i,i)$, $1\le i\le c$, after a translation. The lower boundary of the hexagon between the two
vertical sides consists of the side of length $a$, made of $v$'s, then the side of length $b$, made of $u$'s; so the
midpoints of the right side are $B_j=A_j+av+bu=(a-j,b+j)$. A chain from $A_i$ to $B_j$ has $a+i-j$ east steps among
$a+b$, so $e(A_i,B_j)=\binom{a+b}{a+i-j}$, and disjoint chains are vertex-disjoint paths. Conversely, every family of
vertex-disjoint north-east paths from the $A_i$ to the $B_i$ comes from a tiling. Indeed the hexagon is the set of points
$q+sv+tu+rd$ with $0\le s\le a$, $0\le t\le b$, $0\le r\le c$ (a hexagon with opposite sides parallel is the sum of three
segments), and a path from $A_i$ to $B_i$, which is $A_i$ plus $a$ east and $b$ north steps, stays inside it. The
vertical lozenges along the paths, together with the triangles paired across their vertical edges as in the first part,
then tile the hexagon.

\emph{Determinant.} The sources $A_i=(-i,i)$ move north-west as $i$ increases, and so do the sinks $B_j$. By
Lemma~\ref{linear:lem:planar} non-intersecting systems have $\sigma=\mathrm{id}$, and by Theorem~\ref{linear:thm:lgv} their number is
$\det M$.
\end{proof}

For $a=b=c=2$ the determinant is $\binom42^2-\binom43\binom41=36-16=20$, as announced. The determinant can be evaluated in
closed form, a result that predates the lemma by seventy years.

\begin{remark}[MacMahon's formula]
MacMahon proved in 1896--1916 that the number of plane partitions fitting in an $a\times b\times c$ box, which is the
number of lozenge tilings of $H(a,b,c)$ (Figure~\ref{linear:fig:lozenge} shows why: a tiling is a pile of cubes pushed into a
corner), is
\[
T(a,b,c)=\prod_{i=1}^a\prod_{j=1}^b\prod_{k=1}^c\frac{i+j+k-1}{i+j+k-2}.
\]
Theorem~\ref{linear:thm:lozenge} reduces this to the evaluation of a binomial determinant, which can be done by elementary row
operations; Krattenthaler's survey (see the Sources) gives the details and many relatives. A direct computation verifies the
formula against the determinant for all $a,b,c\le8$. For $a=b=c=3$ it gives $980$. For random tilings of a large regular
hexagon, Cohn, Larsen and Propp proved in 1998 that with probability tending to $1$ the six corners are frozen (tiled by
a single kind of lozenge) outside the circle inscribed in the hexagon, and disordered inside: the \emph{arctic circle}.
The cover of this book shows a sample.
\end{remark}

\begin{beauty}
A determinant is a signed sum of $n!$ products, most of them meaningless. The lemma of Lindstr\"om, Gessel and Viennot
says that for path counts, all the meaningless terms cancel in pairs by a single local move, the swapping of two tails at
a crossing, and what survives is exactly the set of configurations in which nothing crosses. Determinants thus count
non-colliding objects; in probability the same mechanism, discovered by Karlin and McGregor in 1959, gives the
probability that independent random walkers never meet, and it lies behind the eigenvalue statistics of random
matrices (Volume~11).
\end{beauty}

\section{Exercises}

\begin{exercise}[$\star$]\label{linear:ex:latticetriangle}
Compute the area of the polygon with vertices $(0,0),(4,0),(5,3),(2,5),(-1,2)$ with the shoelace formula. Show that a
triangle with vertices in $\Z^2$ has area in $\frac12\Z$.
\end{exercise}

\begin{exercise}[$\star$]\label{linear:ex:laplace}
(Laplace expansion along $k$ columns.) For $A\in M_n(\mathbb K)$ and $J\subseteq[n]$ with $|J|=k$, show that
$\det A=\sum_{|I|=k}(-1)^{s(I)+s(J)}\det A_{I,J}\det A_{I^c,J^c}$. (Use the universal property: for fixed columns outside $J$,
the determinant is an alternating function of the columns in $J$.)
\end{exercise}

\begin{exercise}[$\star\star$]\label{linear:ex:lagrangeid}
(a) Prove Lagrange's identity $\bigl(\sum a_i^2\bigr)\bigl(\sum b_i^2\bigr)-\bigl(\sum a_ib_i\bigr)^2=\sum_{i<j}(a_ib_j-a_jb_i)^2$
with Cauchy--Binet, and deduce the Cauchy--Schwarz inequality with its case of equality. (b) For $v,w\in\R^3$ show that
$\|v\times w\|^2=\|v\|^2\|w\|^2-(v\cdot w)^2$.
\end{exercise}

\begin{exercise}[$\star\star$]\label{linear:ex:charpoly}
Show that the characteristic polynomial of $A\in M_n(\mathbb K)$ is $\det(XI-A)=\sum_{k=0}^n(-1)^k\operatorname{tr}(C_k(A))X^{n-k}$:
the coefficients are, up to sign, the sums of the principal $k\times k$ minors.
\end{exercise}

\begin{exercise}[$\star$]\label{linear:ex:compoundeig}
Show that if $A$ is upper triangular, so is $C_k(A)$ (for the lexicographic order of the subsets), with diagonal entries
$\prod_{i\in I}a_{ii}$. Deduce, using Schur's triangularisation over $\C$ (Appendix~C), that the eigenvalues of $C_k(A)$ are the
products $\lambda_{i_1}\cdots\lambda_{i_k}$, $i_1<\dots<i_k$, of eigenvalues of $A$.
\end{exercise}

\begin{exercise}[$\star\star$]\label{linear:ex:adjadj}
Show that $\operatorname{adj}(\operatorname{adj}A)=(\det A)^{n-2}A$ for $A\in M_n(\mathbb K)$, $n\ge2$.
\end{exercise}

\begin{exercise}[$\star$]\label{linear:ex:dodgson}
Run Dodgson's condensation by hand on $\begin{pmatrix}1&2&3&4\\2&1&2&3\\3&2&1&2\\4&3&2&1\end{pmatrix}$. Then find an invertible
$3\times3$ matrix on which the algorithm divides by $0$, and explain how to repair it.
\end{exercise}

\begin{exercise}[$\star\star$]\label{linear:ex:plucker}
For a $4\times2$ matrix, let $p_{ij}$ be the minor on rows $i<j$. Show that $p_{12}p_{34}-p_{13}p_{24}+p_{14}p_{23}=0$ (a
\emph{Pl\"ucker relation}). Conversely, show that six numbers, not all zero, satisfying this relation are the $2\times2$ minors
of some $4\times2$ matrix of rank $2$. (The lines through the origin of $\R^4$ thus form a quadric in a $5$-dimensional
projective space.)
\end{exercise}

\begin{exercise}[$\star\star$]\label{linear:ex:catalan}
Let $C_m=\frac1{m+1}\binom{2m}m$ be the Catalan numbers, which count the north-east lattice paths from $(0,0)$ to $(m,m)$
that never go above the diagonal. Show that $\det(C_{i+j})_{0\le i,j\le n}=1$ and $\det(C_{i+j+1})_{0\le i,j\le n}=1$
for all $n$.
\end{exercise}

\begin{exercise}[$\star\star$]\label{linear:ex:hex2}
Show that $T(a,b,1)=\binom{a+b}a$ (a direct bijection is easy too), and compute $T(a,b,2)$ in closed form. Check MacMahon's
formula in these two cases.
\end{exercise}

\begin{exercise}[$\star\star$]\label{linear:ex:vandermondelgv}
Let $0\le a_1<\dots<a_k$ and $0\le b_1<\dots<b_k$ be integers. Show that $\det\bigl(\binom{a_i}{b_j}\bigr)\ge0$, with equality
if and only if $a_i<b_i$ for some $i$. (Gessel and Viennot, 1985.)
\end{exercise}

\begin{exercise}[$\star\star$]\label{linear:ex:calisson}
(The problem of the calissons.) Show that in every lozenge tiling of $H(a,b,c)$ the three kinds of lozenges appear $ab$,
$bc$ and $ca$ times. In particular, in a regular hexagon of side $n$ each kind appears $n^2$ times, whatever the tiling.
\end{exercise}

\section{Hints}

\begin{hint}{linear:ex:latticetriangle}
Each term $x_ky_{k+1}-x_{k+1}y_k$ is an integer.
\end{hint}
\begin{hint}{linear:ex:laplace}
The map from the columns in $J$ to $\mathbb K$ is alternating $k$-linear, so it factors through $\Lambda^k\mathbb K^n$; evaluate
it on each $e_I$ by moving rows and columns to the front, as in the proof of Theorem~\ref{linear:thm:jacobi}.
\end{hint}
\begin{hint}{linear:ex:lagrangeid}
Apply Corollary~\ref{linear:cor:pythagoras} to the $n\times2$ matrix with columns $a$ and $b$.
\end{hint}
\begin{hint}{linear:ex:charpoly}
Expand $\det(XI-A)$ by multilinearity in the columns, writing column $j$ as $Xe_j-A_j$.
\end{hint}
\begin{hint}{linear:ex:compoundeig}
$\Lambda^kA$ maps $e_I$ to $\bigwedge_{i\in I}Ae_i$ and $Ae_i$ is a combination of $e_1,\dots,e_i$.
\end{hint}
\begin{hint}{linear:ex:adjadj}
For invertible $A$, $\operatorname{adj}A=\det(A)A^{-1}$. Extend by the polynomial argument of Theorem~\ref{linear:thm:desnanot}.
\end{hint}
\begin{hint}{linear:ex:dodgson}
For the repair, permute rows (or use the polynomial trick: condense $A+tI$ and set $t=0$ at the end).
\end{hint}
\begin{hint}{linear:ex:plucker}
Expand the $4\times4$ determinant of the matrix $(M\mid M)$, which is $0$, along its first two columns
(Exercise~\ref{linear:ex:laplace}). For the converse, if $p_{12}\ne0$ take the rows $(1,0),(0,1)$, $(\cdot,\cdot)$, $(\cdot,\cdot)$.
\end{hint}
\begin{hint}{linear:ex:catalan}
Take $A_i=(-i,-i)$ and $B_j=(j,j)$ in the region below the diagonal $y=x$; non-intersecting systems are forced to be
nested.
\end{hint}
\begin{hint}{linear:ex:hex2}
Compute the $2\times2$ determinant and factor it.
\end{hint}
\begin{hint}{linear:ex:vandermondelgv}
North-east paths from $(0,-a)$ to $(b,-b)$ are counted by $\binom ab$. For the positivity, look for a non-intersecting
system made of ``corners''.
\end{hint}
\begin{hint}{linear:ex:calisson}
Count the vertical lozenges chain by chain, and count the triangles.
\end{hint}

\section{Solutions}

\begin{solution}{linear:ex:latticetriangle}
\idea The formula is a sum of $2\times2$ determinants of integer points.
\writeup The terms $x_ky_{k+1}-x_{k+1}y_k$ for the consecutive pairs are $0\cdot0-4\cdot0=0$, $4\cdot3-5\cdot0=12$,
$5\cdot5-2\cdot3=19$, $2\cdot2-(-1)\cdot5=9$, $(-1)\cdot0-0\cdot2=0$, summing to $40$; the area is $20$ (the vertices are listed
counterclockwise, so the sum is positive). For a lattice triangle, the area is $\frac12|\det(q-p,r-p)|$ with an integer
determinant, so it lies in $\frac12\Z$. (Pick's theorem, Volume~8, Chapter~9, refines this.)
\end{solution}

\begin{solution}{linear:ex:laplace}
\idea Freeze the columns outside $J$; the determinant becomes an alternating $k$-linear function of the columns in $J$,
hence a linear function of their exterior product, whose coordinates are the minors $\det A_{I,J}$.
\writeup Let $\varphi(x_1,\dots,x_k)$ be the determinant of the matrix whose columns in $J$ are $x_1,\dots,x_k$ (in order) and
whose other columns are those of $A$. It is alternating and $k$-linear, so $\varphi=\tilde\varphi\circ\wedge$ and
$\det A=\varphi(A_{j_1},\dots,A_{j_k})=\sum_I\det A_{I,J}\,\tilde\varphi(e_I)$. Now $\tilde\varphi(e_I)$ is the determinant of the matrix
whose columns in $J$ are $e_{i_1},\dots,e_{i_k}$. Move the rows $I$ to the top and the columns $J$ to the left, in order, by
$s(I)-k(k+1)/2$ and $s(J)-k(k+1)/2$ adjacent transpositions: the matrix becomes block triangular
$\begin{pmatrix}I_k&*\\0&A_{I^c,J^c}\end{pmatrix}$. So $\tilde\varphi(e_I)=(-1)^{s(I)+s(J)}\det A_{I^c,J^c}$.
\end{solution}

\begin{solution}{linear:ex:lagrangeid}
\idea Cauchy--Binet for $M^{\mathsf T}M$ with $M=(a\mid b)$ gives a Gram determinant on the left and a sum of squared
$2\times2$ minors on the right.
\writeup (a) $M^{\mathsf T}M=\begin{pmatrix}a\cdot a&a\cdot b\\a\cdot b&b\cdot b\end{pmatrix}$, and the $2\times2$ minors of $M$ are
$a_ib_j-a_jb_i$. Corollary~\ref{linear:cor:pythagoras} is Lagrange's identity. The right side is $\ge0$, which is Cauchy--Schwarz,
with equality if and only if all $2\times2$ minors vanish, that is, $M$ has rank at most $1$: $a$ and $b$ are proportional
(Theorem~\ref{linear:thm:exterior}(c)). (b) For $n=3$ the three minors are, up to sign, the coordinates of $v\times w$.
\end{solution}

\begin{solution}{linear:ex:charpoly}
\idea Multilinearity: each column of $XI-A$ is a sum of two vectors, so the determinant is a sum of $2^n$ determinants.
\writeup Column $j$ is $Xe_j-A_j$. Expanding, $\det(XI-A)=\sum_{S\subseteq[n]}X^{n-|S|}D_S$, where $D_S$ is the determinant with
columns $-A_j$ for $j\in S$ and $e_j$ for $j\notin S$. Expanding $D_S$ along the columns $e_j$, $j\notin S$, leaves the
principal minor: $D_S=\det(-A)_{S,S}=(-1)^{|S|}\det A_{S,S}$. Grouping by $k=|S|$ gives
$\sum_k(-1)^kX^{n-k}\sum_{|S|=k}\det A_{S,S}=\sum_k(-1)^k\operatorname{tr}(C_k(A))X^{n-k}$.
\end{solution}

\begin{solution}{linear:ex:compoundeig}
\idea Exterior powers of triangular maps are triangular, because $Ae_i$ only involves $e_1,\dots,e_i$.
\writeup Write $Ae_i=a_{ii}e_i+w_i$, where $w_i$ is a combination of the $e_l$ with $l<i$. Expand
$\Lambda^kA(e_I)=Ae_{i_1}\wedge\dots\wedge Ae_{i_k}$. The term
taking $a_{i_ri_r}e_{i_r}$ from each factor gives $\prod_{i\in I}a_{ii}e_I$; any other term is $e_{l_1}\wedge\dots\wedge e_{l_k}$ with
$l_r\le i_r$ for all $r$ and some $l_r<i_r$, which is $0$ or $\pm e_L$ with $L\neq I$ and $L$ smaller than $I$ in the
order where $L\preceq I$ means: the elements of $L$ sorted are $\le$ those of $I$ sorted. Listing subsets compatibly
with this partial order (the lexicographic order is such a listing), $C_k(A)$ is upper triangular with diagonal entries
$\prod_{i\in I}a_{ii}$. If $A=PTP^{-1}$ with $T$ triangular, then by functoriality $C_k(A)=C_k(P)\,C_k(T)\,C_k(P)^{-1}$, so
$C_k(A)$ has the eigenvalues of $C_k(T)$.
\end{solution}

\begin{solution}{linear:ex:adjadj}
\idea For invertible matrices the adjugate is the inverse up to a scalar, so the identity is a short computation; the
general case follows by the polynomial argument.
\writeup If $A$ is invertible, $\operatorname{adj}A=\det(A)A^{-1}$, with determinant
\[
\det(A)^n\det(A)^{-1}=\det(A)^{n-1}.
\]
Then
\[
\begin{aligned}\operatorname{adj}(\operatorname{adj}A)&=\det(\operatorname{adj}A)(\operatorname{adj}A)^{-1}\eqbr=\det(A)^{n-1}\cdot\frac{A}{\det A}=\det(A)^{n-2}A.\end{aligned}
\]
Both sides have polynomial entries in the entries of $A$, and they agree on all invertible matrices; as in the proof of
Theorem~\ref{linear:thm:desnanot} (replace $A$ by $A+tI$), they agree everywhere.
\end{solution}

\begin{solution}{linear:ex:dodgson}
\idea Follow Corollary~\ref{linear:cor:dodgson} step by step; the denominators come from the interior of the previous step.
\writeup $A^{(2)}$ is the $3\times3$ matrix of adjacent $2\times2$ minors:
\[
A^{(2)}=\begin{pmatrix}1\cdot1-2\cdot2&2\cdot2-3\cdot1&3\cdot3-4\cdot2\\2\cdot2-1\cdot3&1\cdot1-2\cdot2&2\cdot2-3\cdot1\\3\cdot3-2\cdot4&2\cdot2-1\cdot3&1\cdot1-2\cdot2\end{pmatrix}=\begin{pmatrix}-3&1&1\\1&-3&1\\1&1&-3\end{pmatrix}.
\]
$A^{(3)}_{ij}=(A^{(2)}_{ij}A^{(2)}_{i+1,j+1}-A^{(2)}_{i,j+1}A^{(2)}_{i+1,j})/A_{i+1,j+1}$: the entries are
$(9-1)/1=8$, $(1+3)/2=2$, $(1+3)/2=2$, $(9-1)/1=8$, so $A^{(3)}=\begin{pmatrix}8&2\\2&8\end{pmatrix}$. Finally
$\det A=(8\cdot8-2\cdot2)/A^{(2)}_{22}=60/(-3)=-20$. (Check by any other method: $-20$.) A failing example:
$\begin{pmatrix}1&1&0\\1&0&1\\0&1&1\end{pmatrix}$ has determinant $-2$, but its centre entry is $0$, so the last step divides by
$0$. Repairs: permute rows first (here, exchange the first two rows, which changes the sign), or condense $A+tI$
symbolically and set $t=0$ in the final polynomial; Robbins and Rumsey's analysis shows the algorithm then never fails.
\end{solution}

\begin{solution}{linear:ex:plucker}
\idea A Pl\"ucker relation is a Laplace expansion of a determinant that is obviously zero.
\writeup Let $M$ be $4\times2$ and $N=(M\mid M)$, a $4\times4$ matrix with two equal pairs of columns, so $\det N=0$. Expand
along the first two columns (Exercise~\ref{linear:ex:laplace} with $J=\{1,2\}$): $0=\sum_I(-1)^{s(I)+3}p_Ip_{I^c}$. The six terms
pair up ($I$ and $I^c$ give the same product), and $(-1)^{s(I)+3}$ is $+1$ for $I=\{1,2\},\{3,4\},\{1,4\},\{2,3\}$ and
$-1$ for $\{1,3\},\{2,4\}$; halving, $p_{12}p_{34}-p_{13}p_{24}+p_{14}p_{23}=0$. Conversely, suppose the relation holds and
say $p_{12}\ne0$ (the other cases are the same after renumbering). Rescale so that $p_{12}=1$ and take the matrix with rows
$(1,0)$, $(0,1)$, $(-p_{23},p_{13})$, $(-p_{24},p_{14})$. Its minors are $p_{12}=1$, $p_{13}$, $p_{14}$, $p_{23}$, $p_{24}$ (check:
rows $2,3$ give $0\cdot p_{13}-1\cdot(-p_{23})=p_{23}$), and $p_{34}=-p_{23}p_{14}+p_{13}p_{24}$, which is the given value by the
relation.
\end{solution}

\begin{solution}{linear:ex:catalan}
\idea Hankel determinants $\det(c_{i+j})$ are LGV determinants when $c_{i+j}$ counts paths from a point indexed by $i$ to a
point indexed by $j$. Paths below the diagonal from $(-i,-i)$ to $(j,j)$ have $i+j$ steps of each kind: they are counted
by $C_{i+j}$.
\writeup Work in the acyclic graph of north-east steps between points $(x,y)\in\Z^2$ with $y\le x$. With $A_i=(-i,-i)$ and
$B_j=(j,j)$, $e(A_i,B_j)=C_{i+j}$ (translate the path so that it starts at the origin). First, a non-intersecting system
has $\sigma=\mathrm{id}$. Let $i<j$ and suppose that $P$ goes from $A_i$ to $B_l$ and $Q$ from $A_j$ to $B_m$ with $l>m$.
Use the anti-diagonal coordinate of the proof of Lemma~\ref{linear:lem:planar}. At $s=-2i$, $P$ is at $A_i$, with $x=-i$, while $Q$
is at a point $(x,-2i-x)$ of the region, so $-2i-x\le x$, $x\ge-i$: $\delta(-2i)\le0$. At $s=2m$, $Q$ is at $B_m$, with $x=m$,
while $P$ is at a point with $x\ge m$ for the same reason: $\delta(2m)\ge0$. So the paths meet. Next, there is exactly one
non-intersecting system with $\sigma=\mathrm{id}$, made of the \emph{corners} $K_i$: east from $(-i,-i)$ to $(i,-i)$, then north
to $(i,i)$. By induction, suppose the paths from $A_0,\dots,A_{i-1}$ are $K_0,\dots,K_{i-1}$. The path from $A_i$ must start
east (a north step would leave the region) and leave the row $y=-i$ by a north step from some $(x,-i)$ to $(x,-i+1)$; the
points $(x,-i+1)$ of the region with $x\le i-1$ lie on $K_{i-1}$, so $x\ge i$, and since the path ends at $x=i$, $x=i$. So
the path is $K_i$, which meets no $K_k$, $k<i$ (those use rows above $-i$ and columns west of $i$). Hence the
determinant is $1$. For $\det(C_{i+j+1})$, take $B_j=(j+1,j+1)$; the same arguments apply with the corners $K'_i$ going
east from $(-i,-i)$ to $(i+1,-i)$ and north to $(i+1,i+1)$, and the determinant is again $1$.
\end{solution}

\begin{solution}{linear:ex:hex2}
\idea Small determinants, then factorise.
\writeup For $c=1$, $T(a,b,1)=\binom{a+b}a$: tilings of $H(a,b,1)$ are single paths, and MacMahon's product
$\prod_{i,j}\frac{i+j}{i+j-1}$ telescopes over $j$ to $\prod_i\frac{i+b}i=\binom{a+b}a$. For $c=2$,
\[
\begin{aligned}T(a,b,2)&=\binom{a+b}a^2-\binom{a+b}{a+1}\binom{a+b}{a-1}\\&=\binom{a+b}a^2\left(1-\frac{b\,a}{(a+1)(b+1)}\right)=\binom{a+b}a^2\frac{a+b+1}{(a+1)(b+1)},\end{aligned}
\]
using $\binom{a+b}{a+1}=\binom{a+b}a\frac b{a+1}$ and $\binom{a+b}{a-1}=\binom{a+b}a\frac a{b+1}$. For $a=b=2$ this is
$36\cdot\frac59=20$. MacMahon's product for $c=2$: the $k=1$ factors give $\binom{a+b}a$ as before, and the $k=2$ factors
$\prod_{i,j}\frac{i+j+1}{i+j}$ telescope over $j$ to
\[
\prod_i\frac{i+b+1}{i+1}=\frac{(a+b+1)!}{(b+1)!\,(a+1)!}=\binom{a+b}a\frac{a+b+1}{(a+1)(b+1)}:
\]
the same.
\end{solution}

\begin{solution}{linear:ex:vandermondelgv}
\idea Find points in $\Z^2$ such that the number of north-east paths from $A_i$ to $B_j$ is $\binom{a_i}{b_j}$, then apply
LGV and the planarity lemma: the determinant counts non-intersecting systems, so it is $\ge0$, and it is positive as soon as
we exhibit one.
\writeup Let $A_i=(0,-a_i)$ and $B_j=(b_j,-b_j)$. A north-east path from $A_i$ to $B_j$ has $b_j$ east and $a_i-b_j$ north steps,
so $e(A_i,B_j)=\binom{a_i}{b_j}$ (zero if $a_i<b_j$). As $i$ increases the sources move south and the sinks south-east;
reversing the numbering of both (which multiplies the determinant by $\varepsilon(\rho)^2=1$, $\rho$ the reversal) puts them in the
position of Lemma~\ref{linear:lem:planar}. So the determinant counts the non-intersecting systems with $\sigma=\mathrm{id}$, and is
$\ge0$. If $a_i<b_i$ for some $i$, there is no path from $A_i$ to $B_i$, so no such system, and the determinant is $0$. If
$a_i\ge b_i$ for all $i$, let path $i$ go east from $(0,-a_i)$ to $(b_i,-a_i)$, then north to $(b_i,-b_i)$. Call $R_i$ its horizontal part (on the row $y=-a_i$, with $0\le x\le b_i$) and $C_i$ its vertical part (on the column
$x=b_i$, with $-a_i\le y\le-b_i$). For $i<j$ we have $a_i<a_j$ and $b_i<b_j$: the rows $R_i,R_j$ are at different heights, the
columns $C_i,C_j$ at different abscissas, $R_j$ lies below all of $C_i$ (as $-a_j<-a_i$), and $R_i$ stops at $x=b_i<b_j$, before
reaching $C_j$. So the system is non-intersecting and the determinant is $\ge1$.
\end{solution}

\begin{solution}{linear:ex:calisson}
\idea The chains of the proof of Theorem~\ref{linear:thm:lozenge} give the counts of the two vertical kinds at once: each chain
has the same number of steps of each type. The third kind is what is left.
\writeup Each of the $c$ chains is a path with $a$ east steps and $b$ north steps, and each step is one lozenge, of the
kind determined by the step. So the tiling has $ca$ lozenges of one vertical kind and $cb$ of the other. The hexagon
$H(a,b,c)$ is made of $2(ab+bc+ca)$ triangles (it splits into three parallelograms with sides $a,b$, $b,c$ and $c,a$, each made
of $2\times$ its side product triangles), so a tiling has $ab+bc+ca$ lozenges, and $ab$ of them are of the third kind.
For $a=b=c=n$ each kind appears $n^2$ times. The original proof of David and Tomei is the picture of
Figure~\ref{linear:fig:lozenge}: the three kinds are the faces of the cubes seen from the top, the left and the right, and a pile of
cubes in an $n\times n\times n$ corner shows $n^2$ faces in each direction.
\end{solution}

\section*{Sources and further reading}

\begin{itemize}
\item C.~L.~Dodgson, \emph{Condensation of determinants}, Proceedings of the Royal Society of London 15 (1866), 150--155;
and \emph{An Elementary Treatise on Determinants}, Macmillan, 1867.
\item D.~P.~Robbins and H.~Rumsey, \emph{Determinants and alternating sign matrices}, Advances in Mathematics 62 (1986),
169--184; W.~H.~Mills, D.~P.~Robbins and H.~Rumsey, \emph{Alternating sign matrices and descending plane partitions},
Journal of Combinatorial Theory A 34 (1983), 340--359.
\item D.~M.~Bressoud, \emph{Proofs and Confirmations: The Story of the Alternating Sign Matrix Conjecture}, Cambridge
University Press, 1999. A wonderful account of the whole story, with condensation at its start.
\item D.~Zeilberger, \emph{Proof of the alternating sign matrix conjecture}, Electronic Journal of Combinatorics 3(2)
(1996), R13; G.~Kuperberg, \emph{Another proof of the alternating-sign matrix conjecture}, International Mathematics
Research Notices 1996, 139--150.
\item B.~Lindstr\"om, \emph{On the vector representations of induced matroids}, Bulletin of the London Mathematical
Society 5 (1973), 85--90.
\item I.~Gessel and G.~Viennot, \emph{Binomial determinants, paths, and hook length formulae}, Advances in Mathematics 58
(1985), 300--321.
\item S.~Karlin and J.~McGregor, \emph{Coincidence probabilities}, Pacific Journal of Mathematics 9 (1959), 1141--1164.
\item M.~Aigner and G.~M.~Ziegler, \emph{Proofs from THE BOOK}, 6th ed., Springer, 2018, chapter ``Lattice paths and
determinants''.
\item P.~A.~MacMahon, \emph{Combinatory Analysis}, Volume~2, Cambridge University Press, 1916.
\item C.~Krattenthaler, \emph{Advanced determinant calculus}, S\'eminaire Lotharingien de Combinatoire 42 (1999), B42q.
Section~2 evaluates the binomial determinants behind MacMahon's formula.
\item H.~Cohn, M.~Larsen and J.~Propp, \emph{The shape of a typical boxed plane partition}, New York Journal of Mathematics
4 (1998), 137--165. The arctic circle for lozenge tilings.
\item G.~David and C.~Tomei, \emph{The problem of the calissons}, American Mathematical Monthly 96 (1989), 429--431: in
every lozenge tiling of a regular hexagon, the three kinds of lozenges appear equally often.
\item R.~Bhatia, \emph{Matrix Analysis}, Springer, 1997, Chapter~I, for exterior powers and compound matrices.
\end{itemize}
